% ============================================================
%  ch06 —— §4 定理 1 的证明（原文 p.22–p.26）
% ============================================================
\routenote{p.22--p.26}

\section{定理 1 的证明}

本节证明定理 1。

\subsection{第一步：$\Hk_V(A)$ 是 perverse 的}

我们要证明对每个 $V\in\Rep(\ckG)$，复形 $\Hk_V(A_{\varphi,\chi})[-1]$ 是 perverse 层。为简化记号，只考虑层 $A=A_\varphi$；对 $A_{\varphi,\chi}$ 的证明完全相同，只需在论证中把 $G(0,2)$ 处处换为 $G(1,2)$。

回忆 §2.3 的卷积图：
\[
\begin{tikzcd}[column sep=large, row sep=large]
& \Hecke^{\Aff^1}_{G(0,2)} \arrow[dl, "\mathrm{pr}_1"'] \arrow[dr, "\mathrm{pr}_2"] \arrow[d, "\mathrm{pr}_{\Aff^1}"] & \\
\Bun_{G(0,2)} & \Bun_{G(0,2)}\times\Aff^1 \arrow[l, "\mathrm{proj}_1"'] & \Aff^1
\end{tikzcd}
\]
我们的证明基于几个简单的几何观察。首先回忆上图中映射 $\mathrm{pr}_i$ 是局部平凡纤维化：

\begin{remark}
映射 $\mathrm{pr}_1$ 是局部平凡纤维化，即存在光滑图册 $U\to\Bun_{G(0,2)}$，使得
\[
U\times_{\Bun_{G(0,2)}}\Hecke^{\Aff^1}_{G(0,2)}\cong U\times\mathrm{GR}_{G(0,2)}.
\]
此外，$\mathrm{pr}_2$ 在这个图册上也局部平凡，即
\[
(U\times\Aff^1)\times_{\Bun_{G(0,2)}\times\Aff^1}\Hecke^{\Aff^1}_{G(0,2)}\cong U\times\mathrm{GR}_{G(0,2)}.
\]
\end{remark}

\begin{proof}
为找到满足这些条件的图册 $p:U\to\Bun_{G(0,2)}$，我们只需要 $U$ 上对应于 $p$ 的 $G(0,2)$-丛族在 $U\times\Aff^1$ 上平凡。由命题 1.1，对任意 $U$，这个条件在 $U$ 的 étale 拓扑下局部地成立。
\end{proof}

对 $\gamma\in\Omega$，记 $j_\gamma:I(1)/I(2)\hookrightarrow\Bun_{G(0,2)}$ 为典范嵌入。回忆 $A^\gamma=j_{\gamma,!}\mathrm{AS}_\varphi[\dim(\Bun_{G(0,2)})]$。记 $j'_\gamma:\mathrm{pr}_1^{-1}(I(1)/I(2))\hookrightarrow\Hecke^{\Aff^1}_{G(0,2)}$ 为这个嵌入到 Hecke 叠中的原像。

\begin{remark}
$\mathrm{pr}_2$ 到 $\mathrm{pr}_1^{-1}(I(1)/I(2))$ 的限制
\[
\mathrm{pr}_2:\ \mathrm{pr}_1^{-1}(I(1)/I(2))\to\Bun_{G(0,2)}\times\Aff^1
\]
是仿射的。
\end{remark}

\begin{proof}
由推论 1.3，开子集 $I(1)/I(2)\hookrightarrow\Bun_{G(0,2)}$ 由线丛 $L_i$ 的截面的非零性定义，而这些线丛可定义为忠实表示 $G\to GL_n$ 下 $GL_n$ 的相应线丛的拉回。若 $G=GL_n$，则对任意 $x\in\Aff^1$，这些丛的拉回生成 Picard 群 $\mathrm{Gr}_{GL_n(0,2),x}$，后者是 ind-投射概形（例如 [14, Theorem 7 and 8]）。因此 $I(1)/I(2)$ 在 $\mathrm{GR}_{GL_n(0,2)}$ 中的原像在 $\Aff^1$ 上是仿射的。

对一般的 $G$，ind-概形 $\mathrm{Gr}_{G(0,2),x}$ 通常构造为 $\mathrm{Gr}_{GL_n}$ 对某个合适忠实表示的闭子概形（[32]）。故断言又由 $G=GL_n$ 的情形推出。
\end{proof}

现在可以证明第一步。固定任意 $V\in\Rep(\ckG)$ 及 $\Hecke^{\Aff^1}_{G(0,2)}$ 上相应的 perverse 层 $\mathrm{IC}_V$。我们断言注 4.1 蕴涵
\[
j'_!(\mathrm{pr}_1^*\mathrm{AS}_\varphi\otimes\mathrm{IC}_V)\cong j'_*(\mathrm{pr}_1^*\mathrm{AS}_\varphi\otimes\mathrm{IC}_V).
\]
为看出这点，回忆 $j_!$ 与 $j_*$ 的构造与光滑基变换交换。故为证明所断言同构，可取光滑图册 $p:U\to\Bun_{G(0,2)}$ 使得 $\mathrm{pr}_1$ 的拉回同构于投影 $U\times\mathrm{GR}_G\to U$。但此时断言由同构 $j_!\mathrm{AS}_\varphi=j_*\mathrm{AS}_\varphi$（引理 2.3）推出。

特别地，我们发现 $j'_!(\mathrm{pr}_1^*A\otimes\mathrm{IC}_V)[1]$ 是 perverse 层，且
\[
\Hk_V(A)=(\mathrm{pr}_2\circ j')_!(\mathrm{pr}_1^*(A)\otimes\mathrm{IC}_V)
=(\mathrm{pr}_2\circ j')_*(\mathrm{pr}_1^*(A)\otimes\mathrm{IC}_V).
\]
由注 4.2，映射 $(\mathrm{pr}_2\circ j')$ 是仿射的。因此 $(\mathrm{pr}_2\circ j')_*$ 关于 perverse $t$-结构右正合（[3, §4.1.1]）。于是 $\Hk_V(A)[1]$ 必是 perverse 的。

\subsection{第二步：$A$ 是特征层}

我们已在注 2.7 中指出，$I(1)/I(2)$ 在 $\Bun_{G(0,2)}$ 上的作用扩充为卷积图上的作用。故 $\Hk_V(A)$ 是 $(I(1)/I(2),\varphi)$-等变的。由引理 2.3，这意味着对任意 $\gamma\in\pi_0(\Bun_{G(0,2)})$ 有
\[
\Hk_V(A)|_{\Bun^\gamma_{G(0,2)}\times X}=A^\gamma\boxtimes E^\gamma_V,
\]
其中 $E^\gamma_V[1]$ 是 $X$ 上的 perverse 层。

我们断言层 $E^\gamma_V$ 不依赖于 $\gamma$。为此可假设 $\mathrm{IC}_V$ 只支撑在一个连通分支 $\mathrm{GR}^\gamma_G\subset\mathrm{GR}_G$ 上，因为函子 $\Hk_V$ 同构于由 $\mathrm{IC}_V$ 到各连通分支的限制所定义的函子之直和。

注意对任意 $\gamma'\in\Omega$，Hecke 算子 $\Hk_{\gamma'}$ 与 $\Hk_V$ 交换，因为支撑在 $\PP^1$ 上不同点的 Hecke 算子交换。按定义 $A$ 是算子 $\Hk_{\gamma'}$ 的特征层。于是
\[
A^0\boxtimes E^0_V=\Hk_V(A^{-\gamma})=\Hk_V(\Hk_{-\gamma}(A^0))
=\Hk_{-\gamma}(\Hk_V(A^0))=\Hk_{-\gamma}(A^\gamma\boxtimes E^\gamma_V)
=A^0\boxtimes E^\gamma_V.
\]
这说明层 $E^\gamma_V$ 典范同构于 $E^0_V$，故可去掉指标 $\gamma$。

诸 $\Hk_V$ 与表示的张量积相容，故特别有
\[
\Hk_{V\otimes V'}(A)=\Hk_V\circ\Hk_{V'}(A)=A\boxtimes(E_V\otimes E_{V'}).
\]
这意味着 $E_V\otimes E_{V'}[1]$ 仍是 perverse 层。因此 $E_V$ 必是集中在上同调次数 $0$ 的 perverse 层。故对任意 $V$，复形 $E_V$ 是一个层，且 $E_V=j_{\eta,*}E_{V,\eta}$，其中 $j_\eta:\eta\hookrightarrow\PP^1$ 是一般点的包含。

此外，我们得到取值于 $\eta$ 上局部系统范畴的张量函子 $V\mapsto(E_V)_\eta$。这是刚性张量范畴，故由 [12, Theorem 3.2] 它定义 $\eta$ 上的 $\ckG$-局部系统。我们需要证明 $E_{V,\eta}$ 扩充为 $\PP^1\setminus\{0,\infty\}$ 上的局部系统。

对平凡 Hecke 算子 $\Hk_1$ 有典范同构 $E_1=\mathbb{Q}_\ell$，连同映射 $E_1\to E_V\otimes E_{V^*}\to E_1$，其复合等于乘 $\dim(V)$。我们已知 $E_V=j_{\eta,*}E_{V,\eta}$。故对 $\mathbb{G}_m$ 的任意几何点 $x$，纤维 $E_{V,x}$ 是几何一般纤维 $E_{V,\eta}$ 的子空间，且典范映射 $\mathrm{id}_\eta:\mathbb{Q}_{\ell,\eta}\to(E_V\otimes E_{V^*})_\eta$ 经过 $E_{V,x}\otimes E_{V^*,x}$ 分解。这意味着层 $E_V$ 局部常值，因为 $\mathrm{id}_\eta$ 对应于 $E_{V,\eta}$ 的恒等。

于是诸 $E_V$ 定义从 $\Rep(\ckG)$ 到 $\PP^1\setminus\{0,\infty\}$ 上局部系统范畴 $\Loc(\PP^1\setminus\{0,\infty\})$ 的张量函子。由于这又是刚性张量范畴，它定义 $\PP^1\setminus\{0,\infty\}$ 上的 $\ckG$-局部系统。

\subsection{第三步：驯点处的单值}

如定理 1(2) 的陈述那样，现在假设 $\mathbb{G}=G\times\PP^1$ 是常分裂群。为计算 $\Kl^{\ckG}(\varphi)$ 在 $\{0\}$ 处的单值，我们改述 Bezrukavnikov [5] 的一个论证。他的论证依赖于 Gaitsgory [16] 关于中心层的结果以及 Gabber 的结果——邻近闭链上的单值滤过与权滤过一致 [2]。

为解释这个论证，需回忆 Gaitsgory 的构造（[16]）。他考虑了图
\[
\begin{tikzcd}[column sep=large, row sep=large]
\mathrm{GR}^{\PP^1\setminus\{0,\infty\}}_G \arrow[r, hook, "j_{\mathrm{GR}}"] \arrow[d] & \mathrm{GR}_G \arrow[d] & \mathrm{GR}_{G,0}=\mathrm{Fl}_G \arrow[l, hook'] \arrow[d] \\
\PP^1\setminus\{0,\infty\} \arrow[r, hook, "j"] & \Aff^1 & \{0\} \arrow[l, hook']
\end{tikzcd}
\]
以及诱导的邻近闭链函子 $\Psi:\mathrm{Perv}(\mathrm{GR}^{\PP^1\setminus\{0,\infty\}}_G)\to\mathrm{Perv}_{I_0}(\mathrm{Fl}_G)$。他证明了 $\Psi(\mathrm{IC}_V)$ 上的单值作用是幂幺的。

由于映射 $\mathrm{pr}_1:\Hecke^{\Aff^1}_{G(0,2)}\to\Bun_{G(0,2)}$ 局部同构于与纤维 $\mathrm{GR}_{G(0,2)}$ 之积，我们知道
\[
\Psi(\mathrm{pr}_1^*(A)\otimes\mathrm{IC}_V)=\mathrm{pr}_1^*(A)\otimes\Psi(\mathrm{IC}_V).
\]
特别地，由 Gaitsgory 的结果 [16, Theorem 2] 知这个层上的单值作用是幂幺的。因此（[19, Lemma 5.6]）$\mathrm{pr}_{2,!}(\Psi(\mathrm{pr}_1^*(A)\otimes\mathrm{IC}_V))$ 上的单值作用仍是幂幺的。由于取邻近闭链与真推出交换，单值作用在
\begin{equation}
\mathrm{pr}_{2,!}(\mathrm{pr}_1^*(A)\otimes\Psi(\mathrm{IC}_V))=\mathrm{pr}_{2,!}(\Psi(\mathrm{pr}_1^*(A)\otimes\mathrm{IC}_V))=\Psi(A\boxtimes E_V)=A\boxtimes\Psi(E_V)
\end{equation}
上也是幂幺的。按定义，$\Psi(E_V)$ 是 $E_V$ 在 $0$ 的去心形式邻域上方几何点 $\operatorname{Spec}K^{\mathrm{sep}}_0$ 处的茎，带有作为单值的 $\operatorname{Gal}(K^{\mathrm{sep}}_0/K_0)$-作用。把 (4.3) 限制到平凡 $G(0,2)$-丛，于是得到
\begin{equation}
\Kl^V_{\ckG}|_{\operatorname{Spec}K^{\mathrm{sep}}_0}=R\Gamma_c(\mathrm{pr}_1^*A\otimes\Psi(\mathrm{IC}_V)),
\end{equation}
其中惯性群 $I_0\subset\operatorname{Gal}(K^{\mathrm{sep}}_0/K_0)$ 驯顺且幂幺地作用。

我们须证明单值作用由一个主幂幺元给出。回忆仿射旗流形 $\mathrm{Fl}_G=\mathrm{GR}_{G,0}$ 上 $G(0,2)(k[[t]])=I(0)^{\mathrm{opp}}$-轨道由 Iwahori--Weyl 群 $\widetilde W$ 参数化（[32, Prop.~8.1]）。这些轨道闭包的相交上同调层记为 $\mathrm{IC}_{\tilde w}$。与这些层作卷积定义 Hecke 算子 $\Hk_{\tilde w}(A):=\mathrm{pr}_{2,!}(\mathrm{pr}_1^*(A)\otimes\mathrm{IC}_{\tilde w})$。

由 Görtz 与 Haines 的结果 [19, Corollary 1.2]，层 $\Psi(\mathrm{IC}_V)$ 有一个滤链，其伴随分次层同构于 $\mathrm{IC}_w(i)$（$w\in\widetilde W$，$i\in\mathbb{Z}$），且 $\mathrm{IC}_1(i)$ 的重数等于 $\dim H^{2i}(\mathrm{Gr}_G,\mathrm{IC}_V)$。

现在，对任意长度 $l(\tilde w)>0$ 的 $\tilde w\in\widetilde W$，存在单反射 $s$ 使 $\tilde w=\tilde w's$ 且 $l(\tilde w')<l(\tilde w)$。记 $P_s:=I(0)^{\mathrm{opp}}\cup I(0)^{\mathrm{opp}}sI(0)^{\mathrm{opp}}$ 为由 $I(0)^{\mathrm{opp}}$ 与 $s$ 生成的 parahoric 子群，于是投影 $\mathrm{pr}_s:\mathrm{Fl}\to G((t))/P_s$ 是 $\mathbb{P}^1$-丛。由此看出层 $\mathrm{IC}_{\tilde w}$ 形如 $\mathrm{pr}_s^*(\mathrm{IC}_{\tilde w'})$。因此复形 $\Hk_{\tilde w}(A)$ 在更大的 parahoric 子群 $P_s\supset I(0)^{\mathrm{opp}}$ 下不变。由于 $\Hk_{\tilde w}(A)$ 又是 $(I(1)/I(2),\mathrm{AS}_\varphi)$-等变的，这蕴涵 $\Hk_{\tilde w}(A)=0$。

于是只有长度为 $0$ 的 Hecke 算子非平凡地作用在 $A$ 上。因此，由 Görtz 与 Haines 的结果，$\Psi(\Kl^V_{\ckG})=\Kl^V_{\ckG,\eta}$ 的权滤过的维数由 $\dim H^{2i}(\mathrm{Gr}_G,\mathrm{IC}_V)$ 给出。由于单值滤过与权滤过一致 [2]，单值必按主幂幺元作用。这证明了我们的断言。

\begin{remark}
Görtz 与 Haines 文中所需的关键要素是层 $\Psi(\mathrm{IC}_\mu)$ 的迹函数公式。这个迹由 Iwahori--Hecke 代数中心元素的 Bernstein 公式给出（见 [19, §2.7]）。也可以由这个公式直接推出我们的命题。
\end{remark}
